\documentclass[../../../include/open-logic-section]{subfiles}

\begin{document}

\olfileid[hi]{sth}{replacement}{limofsize}
\olsection{आकार-परिसीमन}

प्रतिस्थापन का ``आंतरिक'' औचित्य देने का शायद सबसे सामान्य प्रयास निम्न
धारणा से आता है:
\begin{enumerate}
	\item[] \limofsize. कोई भी वस्तुएँ समुच्चय बनाती हैं, बशर्ते वे बहुत
	अधिक न हों।
\end{enumerate}
यह सिद्धांत तुरंत प्रतिस्थापन का औचित्य सिद्ध करेगा। आखिर प्रतिस्थापन से बना
कोई समुच्चय उस समुच्चय से बड़ा नहीं हो सकता जिससे वह बना है। सटीक रूप में:
मान लें कि प्रतिस्थापन का उपयोग करके आप समुच्चय
$\funimage{\tau}{A} = \Setabs{\tau(x)}{x \in A}$ बनाते हैं; तब
$\cardle{\funimage{\tau}{A}}{A}$; अतः यदि $A$ के !!{element}s समुच्चय बनाने
के लिए बहुत अधिक नहीं थे, तो उनके प्रतिबिंब भी $\funimage{\tau}{A}$ बनाने
के लिए बहुत अधिक नहीं हैं।

प्रतिस्थापन का औचित्य देने के लिए \limofsize{} का आह्वान करने में स्पष्ट
कठिनाई यह है कि हमने अब तक \limofsize{} जैसा कोई सिद्धांत निर्धारित
\emph{नहीं} किया। इसके अतिरिक्त
\crefrange{sth:story::chap}{sth:z::chap} में समुच्चय की संचयी-पुनरावर्ती
अवधारणा की कथा कहते समय किसी बात ने \limofsize{} की दिशा में \emph{संकेत}
तक नहीं किया। वास्तव में इसीलिए बूलोस ने एक स्थान पर लिखा: ``शायद यह
निष्कर्ष निकाला जा सकता है कि समुच्चय सिद्धांत के `पीछे' कम-से-कम दो विचार
हैं'' \citeyearpar[p.~19]{Boolos1989}। एक ओर समुच्चय की संचयी-पुनरावर्ती
अवधारणा से जुड़े विचार~$\Z$ का औचित्य सिद्ध करने के लिए हैं। दूसरी ओर
\limofsize{} प्रतिस्थापन का औचित्य सिद्ध करने के लिए है।

किंतु समस्या केवल यह नहीं कि अब तक हम \limofsize{} के विषय में \emph{मौन}
रहे हैं। समस्या यह है कि \limofsize{} (अभी सूत्रित रूप में) समुच्चय की
संचयी-पुनरावर्ती धारणा के साथ बहुत असहज बैठता है। आखिर उसमें समुच्चयों के
\emph{चरणों} में बनने के विचार का कोई उल्लेख नहीं।

इन ``दो विचारों'' (अर्थात समुच्चय की संचयी-पुनरावर्ती अवधारणा और
\limofsize) के इतिहास को देखते हुए यह बहुत आश्चर्यजनक नहीं। अंततः ये ``दो
विचार'' समुच्चय-सैद्धांतिक विरोधाभास रोकने की दो काफ़ी भिन्न परियोजनाएँ हैं।
समुच्चय की संचयी-पुनरावर्ती धारणा मोटे तौर पर यह कहकर रसेल का विरोधाभास
रोकती है: \emph{हमें कभी रसेल समुच्चय के अस्तित्व की अपेक्षा नहीं करनी चाहिए
थी, क्योंकि वह किसी चरण पर ``बनता'' नहीं।} इसके विपरीत \limofsize{} मोटे
तौर पर यह कहकर रसेल समुच्चय को बाहर करता है: \emph{हमें कभी रसेल समुच्चय के
अस्तित्व की अपेक्षा नहीं करनी चाहिए थी, क्योंकि वह अस्तित्व रखने के लिए बहुत
बड़ा होता।}

इस रूप में बात रखें तो साफ़ कहें: विरोधाभासों के उत्तर के रूप में
\limofsize{} को \emph{बहुत} अधिक औचित्य चाहिए। उदाहरण के लिए रसेल के
विरोधाभास का यह रूप देखें: \emph{ऐसा कोई पग कुत्ता नहीं जो ठीक उन पग
कुत्तों को सूँघे जो स्वयं को नहीं सूँघते}
(\olref[sth][story][rus]{sec} देखें)। यदि आप पूछें ``ऐसा कोई पग क्यों नहीं
है?'', तो यह कहना अच्छा उत्तर नहीं कि ऐसे पग को बहुत अधिक पगों को सूँघना
पड़ता। फिर रसेल समुच्चय के अनस्तित्व की यह सहज व्याख्या अच्छी क्यों होगी कि
वह अस्तित्व रखने के लिए ``बहुत बड़ा'' होता?

संक्षेप में यदि \limofsize{} की ``सहज'' प्रेरणा आपको कुछ रहस्यमय लगे, तो यह
समझने योग्य है।

\end{document}
